%% evaluation.tex
%%

%% ==================
\section{Evaluation}
\label{sec:Evaluation}
%% ==================

Table~\ref{tab:results} reports mean $\pm$ standard deviation across 5 seeds (42-46).

\begin{table}[htp]
\centering
\caption{Results across 5 seeds (mean $\pm$ std)}
\label{tab:results}
\begin{tabular}{lccc}
\hline
Model & Precision & Recall & F1 \\
\hline
ML Detector (rich context) & 1.000 $\pm$ 0.000 & 1.000 $\pm$ 0.000 & 1.000 $\pm$ 0.000 \\
ML Detector (code-only, baseline gap) & 0.768 $\pm$ 0.045 & 0.708 $\pm$ 0.064 & 0.733 $\pm$ 0.019 \\
Rule-Based Only (reference) & 1.000 $\pm$ 0.000 & 0.483 $\pm$ 0.028 & 0.651 $\pm$ 0.025 \\
Hybrid Detector (improved, code-only) & 1.000 $\pm$ 0.000 & 0.483 $\pm$ 0.028 & 0.651 $\pm$ 0.025 \\
\hline
\end{tabular}
\end{table}

%% ===============================
\subsection{Experiment 1: Reproducing the Precision Collapse}
\label{sec:evalOne}
%% ===============================

With comments present, the ML detector is perfect (1.00/1.00/1.00) -- comment language directly names the risk (``WARNING: ...''), which is expected given how the benchmark is constructed. Without comments, precision falls to 0.768 and recall to 0.708, reproducing the qualitative direction of War et al.'s finding: removing natural-language context degrades a text classifier's reliability, not just its confidence.

%% ===============================
\subsection{Experiment 2: The Hybrid Fix -- and Its Real Cost}
\label{sec:evalTwo}
%% ===============================

The hybrid detector restores precision to a perfect 1.000 in the code-only setting, matching the rich-context ideal. But its recall (0.483) is identical to the rule-based-only reference and \emph{lower} than the plain code-only ML classifier's 0.708. The reason is structural, not a tuning failure: the ML classifier contributes nothing beyond the rule layer at this threshold, because on ``hard'' pattern samples (which by construction carry no code-level signal) its predicted probability never exceeds 0.8 -- there is nothing in the code for it to be confident about.

%% ===============================
\subsection{Experiment 3: The Trade-off Is Tunable, Not Fixed}
\label{sec:evalThree}
%% ===============================

Sweeping the hybrid detector's confidence threshold traces a clear Pareto frontier:

\begin{center}
\begin{tabular}{lccc}
\hline
Threshold & Precision & Recall & F1 \\
\hline
0.5 (= plain ML) & 0.768 & 0.708 & 0.733 \\
0.6 & 0.928 & 0.528 & 0.672 \\
0.7 & 0.998 & 0.485 & 0.652 \\
0.8 & 1.000 & 0.483 & 0.651 \\
0.9 & 1.000 & 0.483 & 0.651 \\
\hline
\end{tabular}
\end{center}

Precision rises monotonically with the threshold and plateaus at 1.000 by 0.8; recall falls from the plain classifier's 0.708 down to the rule-only floor of 0.483 and does not go lower, since the rule layer's contribution is threshold-independent. An operator who cannot tolerate false positives should choose a high threshold; one who needs to catch more of the genuinely ambiguous cases, at the cost of some false positives, should choose a lower one. There is no threshold that achieves both -- the hard pattern's missing information cannot be substituted by a smarter combination rule.

%% ===============================
\subsection{Discussion}
\label{sec:discussion}
%% ===============================

This project reproduces War et al.'s finding, builds a working fix for the specific case that fix can address (enumerable, rule-representable misconfigurations), and is explicit about what it cannot fix: misconfigurations whose safety genuinely depends on information that exists only in a comment cannot be recovered by any classifier or rule combination once that comment is gone. This mirrors the paper's own conclusion that semantic context is not merely helpful but load-bearing -- the honest implication is that comment/documentation quality in IaC repositories is itself a security control, not just a readability nicety.
